\scanpage{25}
considérons
\[
b'=b(s'),\quad c'=c(s')
\]
définis par
\[
\begin{aligned}
s'(u)s'(v)&=-c'uv+s'(\deleted{\unsure{-}}b'uv)=-c'uv+s(b'uv)+\tau b'uv\\
&=(\tau b'\deleted{\unsure{uv}}-c')uv\oplus s(b'uv)\\
\Vert&\\
(s(u)\oplus\tau u)(s(v)\oplus\tau v)
 &=\underbrace{s(u)s(v)}_{-cuv+\deleted{\unsure{b}}s(buv)}+
 \overbrace{(\tau u)sv+\tau(v)su}^{s(2\tau uv)}
 +\deleted{(\tau u)(\tau v)}\inserted{\tau^2uv}\\
 &\deleted{=\unsure{-cuv=\tau^2}}\\
 &=(\tau^2-c)uv\oplus s((b\deleted{\unsure{uv+\tau u v+(\tau v)u}}\inserted{+2\tau})uv)
\end{aligned}
\]
donc
\[
\left\{\begin{aligned}b'&=b+2\tau\\
\tau b'-\unsure{\tau'|c'}&=\tau^2-c
\end{aligned}\right.
\qquad\begin{aligned}
\text{i.e. }c'&=c+\tau b'-\tau^2\\
 &=c+\tau b+2\tau^2-\tau^2\\
 &=c+\tau b+\tau^2
\end{aligned}
\]
\[
\left\{\begin{aligned}b'&=b+2\tau\\
c'&=c+\tau b^{\deleted{\unsure{'}}}+\tau^2 .
\end{aligned}\right. \tag*{(17)}
\]

3. Considérons
\[
\deleted{\Delta}\unsure{\mathscr D|\mathcal D|\mathfrak D}\deleted{=\unsure{\Delta(s)}}=b^2-4c\quad\in\Gamma(S,L^{\otimes2}),\tag*{(18)}
\]
on a \deleted{\unsure{si}}
\[
\begin{aligned}
\mathscr D(s')&=(b+2\tau)^2-4(c+\tau b+\tau^2)\\
 &=b^2+(4\tau^2+4\tau b)-4(c+\tau b+\tau^2)\unsure{=}\\
 &=b^2-4c=\mathscr D(s)
\end{aligned}
\]
donc $\Delta$ est un invariant de $\Omega$, indépendant
du splitting $s$, on l'appelle le \uline{discriminant} du
\unsure{revêtemt|revêtement} quadratique $\Omega$ de $S$, ou encore de l'Algèbre
quadratique $B$, \inserted{notation $\mathscr D(\Omega/S)$.} On va \deleted{voir qu'il est}\inserted{le relier au}
\deleted{égal à $-\delta(N_{\Omega/S})$} discriminant de la forme
quadratique $N_{\Omega/S}$ :
\deleted{En effet, la \unsure{f.} forme}
\[
\deleted{\unsure{\Delta|\operatorname{disc}}(N_{\Omega/S})=-4\mathscr D(\Omega/S)} \tag*{\deleted{(19)}}
\]
\deleted{bilinéaire associée à $N_{\Omega/S}$ est \unsure{telle} $\operatorname{tr}(uv)$}

\deleted{i.e.}
\[
\deleted{N(\zeta+\eta)=N(\zeta)+N(\eta)+\operatorname{tr}(\zeta\eta)}\tag*{\deleted{(19)}}
\]
\deleted{\uline{en effet} on peut choisir un splitting $s$, et on a alors, si $\zeta=\lambda\oplus u$, $\eta=\mu\oplus v$, $\zeta\eta=(\lambda\mu\inserted{-cuv})\oplus(\lambda v+\mu u+buv)$}
\[
\begin{aligned}
&\deleted{N(\zeta\eta)=}\ \deleted{N(\zeta+\eta)=(\lambda+\mu)^2+b\inserted{(\lambda+\mu)}(u+v)+c(u+v)^2}\\
&\deleted{=(\lambda^2+b\lambda u+cu^2)+(\mu^2+\unsure{b|b\mu}v+cv^2)+2cuv+2\lambda\mu}
 \ \deleted{\underbrace{+b(\lambda v+\mu u)}}\\
&\deleted{=N(\zeta)+N(\eta)+2cuv+2\lambda\mu+b(\lambda v+\mu u)}
\end{aligned}
\]

% archive sheet number: 24

